\section{Why the proof fits inside iPRA}\label{sec:ipra-audit}

Primitive recursiveness of an algorithm alone would not prove its claimed
correctness in $\ipra$. The relevant point is that the correctness
arguments above use particular finite induction invariants available
in that theory. We record them and the translation into natural-number
arithmetic.

\subsection{The exact first-order assertion}

Fix all evaluator and modulus programs, allowing an optional tuple of
number parameters $a$. Each rational expression $V(a,r)$ or $B(a,r)$
has a primitive recursive code. Write $\mathsf H(a)$ for the universal
arithmetic certificates in \cref{sec:certificates}.
The main precision assertion is the first-order formula
\begin{equation}\label{eq:object-statement}
 \mathsf H(a)\ \longrightarrow\
 \forall r,u,v\,
 \left(r\le u\wedge r\le v\ \longrightarrow\
       64\,2^r\,|V(a,u)-B(a,v)|\le19\right).
\end{equation}
The displayed rational inequality is expanded into a natural-number
comparison after multiplication by positive denominators. The corresponding
formulas for regularity and equality are expanded in the same way.

There is no internal quantifier over real-valued functions in
\eqref{eq:object-statement}. There is no assertion that a single natural
number codes and evaluates every possible primitive recursive program.
For each fixed syntactic instance, the newly constructed names are
themselves fixed primitive recursive definitions. This is the usual
schematic use of the full primitive recursive language.

\subsection{Clearing denominators and handling negative values}

For one requested finite grid, list the rational quantities needed in
the proof: coefficient samples, width-weighted derivative samples,
grid weights, and error terms. Let $D$ be the product of their positive
denominators. There are primitive recursive bounds on the length of
this list, and its construction is a bounded loop over the grid.
Thus $D$ is a positive primitive recursive integer.

After multiplication by $D$, a coordinate line has integer potential
samples $A(j)$, integer derivative increments $b(j)$, and a nonnegative
integer error $E$. Its local contract has the two forms
\[
 A(j+1)-A(j)-b(j)\le E,\qquad
 b(j)-A(j+1)+A(j)\le E.
\]
To use a natural-valued finite-FTC module, set
\[
 L=\max_{j\le N}\max\{0,-A(j)\},\quad
 G(j)=A(j)+L,\quad b(j)=b^+(j)-b^-(j)
\]
on the guarded range $j\le N$. One may totalize $G$ outside that range
using $\max\{0,A(j)+L\}$.
Then $G,b^+,b^-,E$ are natural-valued primitive recursive functions
on the needed indices. The first local inequality becomes
\[
 G(j+1)+b^-(j)\le G(j)+b^+(j)+E,
\]
and the second has the sides reversed appropriately.
The same offset $L$ occurs at both endpoints and cancels.
No positivity hypothesis on the original coefficient is introduced.

If multiple lines are processed together, take one denominator and one
offset for the entire finite request, or package per-line values into
the extra parameters. Both choices are primitive recursive. The
common-denominator variant is slower but particularly simple to
formalize. Reduction of fractions is unnecessary.

\subsection{The induction invariants}

The following ledger identifies the actual objects of induction.

\begin{longtable}{>{\raggedright\arraybackslash}p{0.22\textwidth}p{0.70\textwidth}}
\toprule
Argument & Quantifier-free invariant after coding\\
\midrule
\endhead
Telescoping &
The denominator-cleared equality
$\sum_{j<m}(A(j+1)-A(j))=A(m)-A(0)$.
The successor case adds one difference and cancels one term.\\
\addlinespace
Error accumulation &
The two inequalities bounding the signed prefix residual by $mE$,
guarded by $m\le N$. The successor case adds the local inequality to
the previous prefix inequality.\\
\addlinespace
Finite Fubini &
Equality of two nested finite sums on a rectangular prefix.
At the successor, distribute the newly added row and use the inner
finite sum identity.\\
\addlinespace
Refinement &
The partial grouped-sum identity and its accumulated
$\sum w_j\eta$ error bound, with a denominator-cleared weight sum.
One additional fine point contributes one (U) instance.\\
\addlinespace
Boolean finite checks &
A prefix fold of decidable rational certificate tests equals one.
This is an equation between primitive recursive terms.\\
\addlinespace
Exterior algebra &
Coefficient equalities for finite product and permutation expansions.
The invariants are finite polynomial equalities and inversion-count
parity identities.\\
\addlinespace
Double boundary &
The sum of paired codimension-two coefficients is zero.
Each finite pair consists of equal terms with opposite signs.\\
\bottomrule
\end{longtable}

The universal hypotheses (C), (U), and (D) stay in the proof context.
Their specific rational instances provide the local inequalities.
They are not inserted into an unbounded analytic induction invariant.
When a bounded universal condition is convenient in a finite invariant,
replace it by its Boolean fold. Its equivalence to the bounded
condition follows by a finite induction because each atomic test is
decidable.

The final Cauchy and equality statements are derived by introducing
arbitrary precision indices and applying the finite lemmas. There is no
induction on a formula asserting convergence at every precision.
The order split $u\le v$ or $v\le u$ is a decidable comparison of
natural numbers. The rational slack lemma similarly uses decidable
rational order and a specifically computed precision, not a principle
for deciding the order of represented reals.

\subsection{The remaining calculus and geometry}

The product and chain rules use fixed finite algebraic expansions,
integer bounds, and maximum/precision formulas. A purported application
to a particular object must supply the jet hypotheses, but their
closure does not require a stronger induction principle.
The use of named arguments is interpreted by
\eqref{eq:extension}; the only removal of auxiliary errors is
\cref{lem:slack}, whose witness is primitive recursive.

Each exterior coefficient is a finite expression built from those
operations. Naturality is obtained by finite expansion, product rules,
and the supplied symmetric second jets. The chain step is a finite
integer linear combination with the displayed weight shift.
These steps therefore extend the cube derivation schematically in the
same first-order theory.

This supplies a human-readable proof of the iPRA derivability claim.
A machine-checked realization must still instantiate these invariant
patterns in the chosen syntax and prove the necessary substitution
lemmas. The following section gives that implementation plan.

\section{Lean formalization and reuse of ProveIt}\label{sec:lean}

\subsection{Three distinct validation layers}

The implementation should keep the following claims separate.

\begin{enumerate}
\item \textbf{Host finite mathematics.} A Lean theorem about integer or
rational sums checks the finite identity in Lean's ambient logic.
\item \textbf{Object-level derivability.} A Lean term of type
$\Gamma\vdash_{\mathsf{PRAAx}}\varphi$ is a derivation of the encoded
formula using the specified iPRA axioms and rules.
\item \textbf{Effective analytic interpretation.} The encoded formula
must be proved to express the advertised evaluations, moduli,
integrals, and equality relation. This includes the signed-rational
adapter and the passage from finite grids to names.
\end{enumerate}

The existing Mathlib divergence theorems are useful mathematical
comparison points, but their statements use real normed spaces and
Henstock--Kurzweil or Bochner integration
\cite{mathlib-box,mathlib-bochner}. Invoking one of them would not
establish the object-level iPRA theorem. Our intended formalization
instead follows the finite rational argument above.

\subsection{Inspected repository lemmas}

All paths in the table are relative to
\path{Logic/PrimitiveRecursiveArithmetic/Derivations/Lean/IPRADerivations/}.
The source snapshot is \RepoCommit.

\begin{longtable}{>{\raggedright\arraybackslash}p{0.35\textwidth}p{0.56\textwidth}}
\toprule
Module and names & Actual role in the planned proof\\
\midrule
\endhead
\path{Analysis/FiniteFTC.lean}\newline
\texttt{ftc\_upper}, \texttt{ftc\_lower}
& Object derivations of the two accumulated linewise inequalities.
The source explicitly uses guarded quantifier-free induction.\\
\addlinespace
\path{Core/Fubini.lean}\newline
\texttt{LeanProofs.IPRA.fubini}
& Object-level exchange of two bounded sums of a fixed primitive
recursive body. Repeated exchanges move one chosen coordinate
innermost.\\
\addlinespace
\path{Core/SumLemmas.lean}\newline
\texttt{sum\_add}, \texttt{sum\_shift}, \texttt{sum\_congr\_bounded}
& Finite sum arithmetic and replacement of equal bounded summands.\\
\addlinespace
\path{Analysis/RiemannSumRules.lean}\newline
\texttt{rs\_close}, \texttt{rs\_close\_const}, \texttt{rs\_bound}
& Object derivations that sum pointwise inequalities and accumulate
uniform error budgets.\\
\addlinespace
\path{Analysis/RiemannRefinement.lean}\newline
\texttt{refine\_upper}, \texttt{refine\_lower}
& An inspected-source Lipschitz refinement pattern.
The arbitrary supplied-modulus version in this report is an
additional obligation.\\
\addlinespace
\path{Analysis/TaggedFTC.lean}\newline
\texttt{tagged\_upper}, \texttt{tagged\_lower}
& Variable-width finite estimates with evaluation-error accumulation.
The source does not itself derive every analytic jet input or
the passage to integral names.\\
\addlinespace
\path{Foundations/DyadicNames.lean}\newline
\texttt{sample\_stability}, \texttt{add\_regular}
& Arithmetic on dyadic errors after common scaling, with name
addition estimates.\\
\addlinespace
\path{Foundations/UniformCertificates.lean}\newline
\texttt{grid\_clip}, \texttt{net\_error}, \texttt{error\_transport}
& Finite-grid and error-transport certificates. These do not
extract moduli from arbitrary pointwise continuity.\\
\bottomrule
\end{longtable}

In particular, the upper finite-FTC theorem accepts a derivation of
\[
 \forall j<N,\quad G(j+1)+F^-(j)\le G(j)+F^+(j)+E
\]
and returns a derivation of
\[
 G(N)+\sum_{j<N}F^-(j)
 \le G(0)+\sum_{j<N}F^+(j)+NE.
\]
This is exactly the natural-number form of the line estimate described
in \cref{sec:ipra-audit}. The lower theorem supplies the other direction.
The rational and multidimensional adapters are genuine work still to
be implemented, but the finite induction kernel is already concretely
identified.

Two nearby files must not be overinterpreted.
\path{Topology/BoundarySquare.lean} contains an ambient Lean simplicial
boundary calculation, not already the internal cubical theorem.
\path{Analysis/PullbackCertificates.lean} concerns hypergeometric
substitutions, not differential-form pullbacks. The source manifest
records the exact inspected paths, hashes, imports, and theorem names.

\subsection{Definitions to implement}

The accompanying \path{lean/StokesBlueprint.lean.txt} specifies the
contracts and intended statements. A practical module decomposition is:

\begin{enumerate}
\item \textbf{Rational codes and primitive recursive terms.}
Define signed pairs, positive denominators, cross-multiplied
comparison, powers, bounded sums and maxima, the coded sup norm,
and their substitution rules.
\item \textbf{Grid terms.}
Define base-$N$ decoding of the $N^k$ grid indices, insertion of a
coordinate, upper/lower face maps, and the terms $V,B,T$.
Prove cardinality and weight identities.
\item \textbf{Local contracts.}
Encode (C), (U), and (D) as first-order formulas in the fixed terms.
The finite local residual bound is the direct input to the signed
adapter.
\item \textbf{Finite Stokes.}
Build the common denominator and offset; invoke
\texttt{ftc\_upper}/\texttt{ftc\_lower} on each line;
apply finite Fubini and \texttt{rs\_close};
derive \eqref{eq:finite-bound}.
\item \textbf{Names and integration.}
Prove the arbitrary-modulus refinement lemma, the regular integral
construction, evaluator independence, rational slack elimination,
and the schedule inequalities with constants $7/64,3/32,3/16$.
\item \textbf{Jets and forms.}
Implement the finite-expression evaluator, clipped named-input
composition, derivative contracts, finite exterior words and
minors, and the coefficientwise naturality proof.
\item \textbf{Chains and surface corollaries.}
Represent frozen-coordinate face maps, prove their normalization
and sign pairing, add supplied face-matching witnesses, and
combine the finite precision budgets.
\end{enumerate}

Each item has its mathematical definition and proof in this report.
The remaining task is to construct the actual syntax-level derivations
and run the compiler, rather than to find an additional analytic
principle.

\subsection{The supplied Lean files}

\path{lean/FiniteCore.lean} contains standard-library proof scripts for
finite integer telescoping, upper/lower error accumulation, transverse
row cancellation, and square double-boundary cancellation. These
are host-level theorems.

\path{lean/RepositorySlice.lean} targets the actual ProveIt derivation
judgments. It exposes the two finite-FTC adapters and supplies
polynomial certificates for the square-boundary identity and the
three-dimensional curl/pullback algebra. Its polynomial identity
does not by itself prove the analytic product or chain rule.

\path{lean/StokesBlueprint.lean.txt} is intentionally a blueprint,
not a compilable module. It gives the complete intended interface,
proof invariants, coding choices, and order of implementation.
The two \path{.lean} scripts contain no intentional proof-hole or
new-axiom declarations, but \emph{they were not compiled in this
environment}. Absence of such declarations is not a substitute
for compilation.

The inspected repository inventory specifies Lean \texttt{v4.32.0}
and Mathlib commit
\begin{center}
\small\texttt{81a5d257c8e410db227a6665ed08f64fea08e997}.
\end{center}
These are source metadata, not a fresh build receipt. The adapter
requires the full pinned repository and its canonical dependency
configuration; the archive does not purport to contain that
repository's entire transitive dependency closure.

\subsection{Reproduction and completion gates}

The executable finite checks are run from the package root by
\begin{lstlisting}
python3 code/validate.py
\end{lstlisting}
The optional figures are regenerated with
\begin{lstlisting}
python3 code/make_figures.py
\end{lstlisting}
The report source can be compiled by a standard LaTeX installation;
\path{build_report.sh} also supports the installed-file fallback used
to build this delivery. With Lean installed, first check the
standalone core, then the adapter in the complete pinned checkout.
The precise replay instructions are in \path{lean/README.md}.

For a completed formalization the final acceptance condition is not
merely a successful theorem about $\R$. It is a clean build of the
specific object-level derivation of \eqref{eq:object-statement},
followed by its named-integral interpretation, with the proof
dependencies and axioms inspected. The present delivery supplies
the complete human proof, an explicit route to that result, and
executable finite checks; it does not claim this final machine
verification has already happened.

\section{Scope boundaries and useful extensions}\label{sec:boundaries}

The theorem proves the integral identity for supplied effective cell
presentations; uniform moduli and a finite chain are input data.
The $C^2$ parameterization interface suffices for the elementary
pullback argument. Weaker geometric regularity, singular boundaries,
or unbounded domains would require additional explicit estimates.

Rational boxes via affine pullback, dyadic subdivision, coefficients
in fixed finite-dimensional rational normed spaces, and comparisons
with supplied fillings fit the same finite method. Automatically
producing uniform moduli or finite geometric presentations remains
a separate task.
